\documentclass[11pt,a4paper]{article}

\usepackage{amsmath,amssymb,amsthm}
\usepackage{bm}
\usepackage{geometry}
\usepackage{booktabs}
\usepackage{xcolor}

\geometry{margin=2.5cm}

\theoremstyle{plain}
\newtheorem{theorem}{Theorem}
\newtheorem{corollary}[theorem]{Corollary}
\newtheorem{proposition}[theorem]{Proposition}
\theoremstyle{definition}
\newtheorem{definition}{Definition}

\newcommand{\freq}{\tilde{f}}
\newcommand{\fstar}{f^*}
\newcommand{\fhat}{\hat{f}}

\title{Mathematical Relation Between True Subgraph Frequency\\and OES-Estimated Frequency in Minomaly}
\author{}
\date{}

\begin{document}
\maketitle

\section{Setup and Notation}

Let $G = (V, E)$ be the host graph with $n = |V|$ nodes. Let $\mathcal{G}(G)$ denote the set of all connected subgraphs of~$G$, and let $\mathcal{D}(G)$ be the probability distribution over $\mathcal{G}(G)$ induced by the tree sampling procedure (BFS frontier walk with anchor node, random size in $[s_{\min}, s_{\max}]$).

Let $RG = \{G'_1, G'_2, \ldots, G'_k\}$ be $k$ subgraphs sampled i.i.d.\ from $\mathcal{D}(G)$, and let $\phi : \mathcal{G}(G) \to \mathbb{R}^d$ be the learned order embedding function.

We denote subgraph isomorphism by $P \subseteq_i H$, meaning $\exists\, H' \subseteq H$ such that $H' \cong P$.

\begin{definition}[True Frequency]
For a pattern $P \in \mathcal{G}(G)$, the \textbf{true frequency} under the sampling distribution is:
\[
\fstar(P) \;=\; \Pr_{H \sim \mathcal{D}(G)}\!\big[P \subseteq_i H\big]
\]
Computing $\fstar$ requires solving the NP-complete subgraph isomorphism problem for every pair $(P, H)$.
\end{definition}

\begin{definition}[Sample Frequency]
Given a perfect subgraph-isomorphism oracle, the \textbf{sample frequency} on $RG$ is:
\[
\fhat(P) \;=\; \frac{1}{k}\sum_{i=1}^{k} \mathbf{1}\big[P \subseteq_i G'_i\big]
\]
\end{definition}

\begin{definition}[OES-Estimated Frequency]
The \textbf{OES-estimated frequency} computed by Minomaly is:
\[
\freq(P) \;=\; \frac{1}{k}\sum_{i=1}^{k} \mathbf{1}\big[\phi(P) \preceq \phi(G'_i)\big]
\]
where $\phi(P) \preceq \phi(G'_i)$ denotes the component-wise partial order as predicted by the classifier: $\text{clf\_model}\!\big(E(P, G'_i)\big) = 1$, with the order-embedding violation
\[
E(P, H) = \sum_{j=1}^{d} \big[\max\!\big(0,\; \phi_j(P) - \phi_j(H)\big)\big]^2.
\]
\end{definition}

\section{Error Sources}

There are exactly two sources of discrepancy between $\freq(P)$ and $\fstar(P)$.

\subsection{Sampling Error}

Since $RG$ is a finite sample from $\mathcal{D}(G)$, the sample frequency $\fhat(P)$ deviates from $\fstar(P)$ by random noise. Each $\mathbf{1}[P \subseteq_i G'_i]$ is a Bernoulli random variable with parameter $\fstar(P)$, so $\fhat(P)$ is an \emph{unbiased} estimator:
\[
\mathbb{E}\big[\fhat(P)\big] = \fstar(P).
\]

\subsection{Embedding Error}

The order embedding is imperfect. We define the error rates:

\begin{definition}[Embedding Error Rates]
\begin{align*}
\alpha &= \Pr\!\big[\phi(P) \preceq \phi(H) \;\big|\; P \not\subseteq_i H\big] && \text{(false positive rate)} \\[4pt]
\beta  &= \Pr\!\big[\phi(P) \not\preceq \phi(H) \;\big|\; P \subseteq_i H\big] && \text{(false negative rate)}
\end{align*}
We define the \textbf{discriminability} as $\gamma = 1 - \alpha - \beta$.
\end{definition}

\noindent The reported 95\% accuracy of the pretrained model on balanced positive/negative pairs implies $\frac{(1-\alpha)+(1-\beta)}{2} \approx 0.95$, i.e., $\alpha + \beta \approx 0.05$ and $\gamma \approx 0.95$.

\section{Main Result}

\begin{theorem}[Linear Bias-Noise Decomposition]\label{thm:main}
Let $\alpha, \beta$ be the false positive and false negative rates of the order embedding, assumed uniform across patterns. Let $\gamma = 1 - \alpha - \beta$. Then:
\[
\boxed{\;\freq(P) \;=\; \gamma \cdot \fstar(P) \;+\; \alpha \;+\; \varepsilon(P, k)\;}
\]
where the error term $\varepsilon(P, k)$ satisfies
\[
\big|\varepsilon(P, k)\big| \;\leq\; \sqrt{\frac{\ln(2/\delta)}{2k}}
\]
with probability at least $1 - \delta$.
\end{theorem}

\begin{proof}
For each $G'_i \in RG$, the indicator $X_i = \mathbf{1}[\phi(P) \preceq \phi(G'_i)]$ is a Bernoulli random variable. By the law of total probability, its parameter is:
\begin{align*}
q_i &:= \Pr\!\big[\phi(P) \preceq \phi(G'_i)\big] \\
    &= (1-\beta)\cdot\mathbf{1}[P \subseteq_i G'_i] \;+\; \alpha\cdot\mathbf{1}[P \not\subseteq_i G'_i].
\end{align*}

Since $G'_i \sim \mathcal{D}(G)$ independently, taking the expectation over the sampling distribution:
\begin{align*}
\mathbb{E}[q_i] &= (1-\beta)\cdot\Pr[P \subseteq_i G'_i] + \alpha\cdot\Pr[P \not\subseteq_i G'_i] \\
                 &= (1-\beta)\cdot \fstar(P) + \alpha\cdot\big(1 - \fstar(P)\big) \\
                 &= \fstar(P) - \beta\,\fstar(P) + \alpha - \alpha\,\fstar(P) \\
                 &= (1 - \alpha - \beta)\,\fstar(P) + \alpha \\
                 &= \gamma\cdot \fstar(P) + \alpha.
\end{align*}

Since $\freq(P) = \frac{1}{k}\sum_{i=1}^{k} X_i$ is the mean of $k$ independent random variables with $X_i \in [0,1]$, Hoeffding's inequality gives:
\[
\Pr\!\Big[\big|\freq(P) - \mathbb{E}[\freq(P)]\big| \geq t\Big] \;\leq\; 2\exp(-2kt^2).
\]

Setting $t = \sqrt{\ln(2/\delta)/(2k)}$ and defining $\varepsilon(P,k) = \freq(P) - \mathbb{E}[\freq(P)]$ completes the proof.
\end{proof}

\section{Corollaries}

\begin{corollary}[Corrected Frequency Estimator]\label{cor:corrected}
An approximately unbiased estimator of $\fstar(P)$ is:
\[
\fstar_{\mathrm{corr}}(P) \;=\; \frac{\freq(P) - \alpha}{\gamma}
\]
with error bound:
\[
\big|\fstar_{\mathrm{corr}}(P) - \fstar(P)\big| \;\leq\; \frac{1}{\gamma}\sqrt{\frac{\ln(2/\delta)}{2k}}
\]
with probability at least $1 - \delta$.
\end{corollary}

\begin{proof}
Immediate from Theorem~\ref{thm:main}: $\fstar_{\mathrm{corr}}(P) = \frac{\freq(P) - \alpha}{\gamma} = \fstar(P) + \frac{\varepsilon(P,k)}{\gamma}$.
\end{proof}

\noindent\textbf{Numerical example.} For $\gamma = 0.95$, $k = 10{,}000$, $\delta = 0.05$:
\[
\big|\fstar_{\mathrm{corr}}(P) - \fstar(P)\big| \;\leq\; \frac{1}{0.95}\sqrt{\frac{\ln 40}{20{,}000}} \;\approx\; 0.0143.
\]
The corrected OES frequency approximates the true frequency \textbf{within 1.4\%} at 95\% confidence.

\begin{corollary}[Effective Anomaly Threshold]\label{cor:threshold}
When Minomaly flags a pattern as anomalous at threshold $T_F$ (i.e., $\freq(P) < T_F/n$), the effective condition on the true frequency is:
\[
\fstar(P) \;<\; \frac{T_F/n - \alpha}{\gamma} + O(k^{-1/2}).
\]
\end{corollary}

\begin{proof}
From Theorem~\ref{thm:main}, $\freq(P) < T_F/n$ implies $\gamma\cdot\fstar(P) + \alpha + \varepsilon < T_F/n$. Rearranging and bounding $|\varepsilon| = O(k^{-1/2})$ gives the result.
\end{proof}

\begin{corollary}[Preservation of Anti-Monotonicity]\label{cor:antimon}
Let $P_1 \subseteq P_2$, so that $\fstar(P_1) \geq \fstar(P_2)$ by the anti-monotonicity of true frequency. If the frequency gap satisfies:
\[
\fstar(P_1) - \fstar(P_2) \;>\; \frac{2}{\gamma}\sqrt{\frac{\ln(4/\delta)}{2k}},
\]
then $\freq(P_1) > \freq(P_2)$ with probability at least $1 - \delta$.
\end{corollary}

\begin{proof}
By Theorem~\ref{thm:main} applied to each pattern:
\begin{align*}
\freq(P_1) &\geq \gamma\,\fstar(P_1) + \alpha - |\varepsilon_1|, \\
\freq(P_2) &\leq \gamma\,\fstar(P_2) + \alpha + |\varepsilon_2|.
\end{align*}
Subtracting:
\[
\freq(P_1) - \freq(P_2) \;\geq\; \gamma\big(\fstar(P_1) - \fstar(P_2)\big) - |\varepsilon_1| - |\varepsilon_2|.
\]
By the union bound, both $|\varepsilon_i| \leq \sqrt{\ln(4/\delta)/(2k)}$ hold simultaneously with probability $\geq 1 - \delta$. Substituting the gap condition gives $\freq(P_1) - \freq(P_2) > 0$.
\end{proof}

\begin{corollary}[Minimum Sample Size]\label{cor:kmin}
To guarantee that $\freq$ preserves the frequency ordering of two patterns separated by a true frequency gap~$\Delta$, with confidence $1 - \delta$, the required sample size is:
\[
k \;\geq\; \frac{2\ln(4/\delta)}{\gamma^2\,\Delta^2}.
\]
\end{corollary}

\begin{proof}
Rearranging the condition in Corollary~\ref{cor:antimon}: $\Delta > \frac{2}{\gamma}\sqrt{\frac{\ln(4/\delta)}{2k}}$ is equivalent to $k > \frac{2\ln(4/\delta)}{\gamma^2 \Delta^2}$.
\end{proof}

\noindent\textbf{Numerical examples.}

\begin{table}[h]
\centering
\caption{Minimum $k$ for anti-monotonicity preservation ($\gamma=0.95$, $\delta=0.05$)}
\label{tab:kmin}
\begin{tabular}{lccc}
\toprule
\textbf{Frequency gap} $\Delta$ & \textbf{Min $k$} & \textbf{Cora} ($k\!=\!10\text{K}$) & \textbf{Flickr} ($k\!=\!250\text{K}$) \\
\midrule
$5\%$  & $3{,}247$    & \checkmark & \checkmark \\
$3\%$  & $9{,}020$    & \checkmark & \checkmark \\
$1\%$  & $81{,}178$   & $\times$   & \checkmark \\
$0.5\%$ & $324{,}711$ & $\times$   & $\times$    \\
\bottomrule
\end{tabular}
\end{table}

\noindent This explains the experimental parameter choices: Cora ($k = 10{,}000$) suffices for frequency gaps $\geq 3\%$, while Flickr ($k = 150{,}000$--$250{,}000$) is needed for finer discrimination.

\section{Interpretation}

The central equation from Theorem~\ref{thm:main},
\[
\freq(P) = \gamma \cdot \fstar(P) + \alpha + \varepsilon(P,k),
\]
reveals that the OES-estimated frequency is a \textbf{linearly biased, noisy estimator} of the true frequency:

\begin{itemize}
    \item \textbf{Slope} $\gamma = 1 - \alpha - \beta \approx 0.95$: the discriminability compresses the frequency range. A pattern with true frequency 50\% is estimated at $\approx 0.95 \times 0.50 + 0.025 = 50\%$, while a pattern with true frequency 1\% is estimated at $\approx 0.95 \times 0.01 + 0.025 = 3.5\%$. The compression slightly reduces contrast between frequent and infrequent patterns.

    \item \textbf{Intercept} $\alpha \approx 0.025$: false positives shift all estimated frequencies upward. Even a pattern that truly appears nowhere ($\fstar = 0$) would be estimated at $\freq \approx \alpha = 2.5\%$. This sets a \emph{noise floor}: patterns with true frequency below $\alpha$ cannot be reliably distinguished from zero-frequency patterns.

    \item \textbf{Noise} $\varepsilon = O(k^{-1/2})$: sampling variance that vanishes with larger~$k$. This is the standard Monte Carlo convergence rate, unaffected by the graph size~$n$.
\end{itemize}

\paragraph{Implications for anomaly detection.}
The linear relation means that the \emph{ranking} of patterns by frequency is preserved (up to the noise level), even though the absolute values are biased. Since Minomaly only needs the ranking to be correct (to identify the rarest patterns and guide the greedy search), the bias $\alpha$ does not affect detection quality as long as $T_F/n > \alpha$. The noise floor $\alpha$ effectively sets the minimum detectable frequency: patterns rarer than $\alpha$ in truth cannot be distinguished from each other.

\paragraph{Relation to the clf\_model.}
In practice, the false positive rate $\alpha$ is not a fixed global constant --- the learned classifier (a Linear$(1, 2)$ + LogSoftmax layer applied to the scalar violation $E$) adapts its decision threshold to the observed distribution of violations. For well-trained models, $\alpha$ is typically much smaller than the balanced accuracy would suggest, because the violation distributions of true positives ($E \approx 0$) and true negatives ($E \gg \alpha_{\text{margin}}$) are well-separated. The 95\% accuracy reported on balanced synthetic test sets represents a conservative upper bound on $\alpha + \beta$ for the in-distribution comparisons encountered during detection.

\end{document}
